\scanpage{204}\mspage{149}
\deleted{\unsure{donc}} \inserted{donc} un foncteur, \unsure{définit|défini} à isom unique près :
\[
 f_{\unsure{0|\circ}}^*:\unsure{A|\mathcal{A}}_0\longrightarrow A'_0,
\]
donnant lieu à un carré de foncteurs, commutatif à isom fonctoriel canonique près
\diagram{\begin{tikzcd}[ampersand replacement=\&]
 \& A_0 \arrow[r,"i"] \arrow[d,"f_0^*"'] \& A \arrow[d,"f^*"]\\
 (\mathrm{II}) \& A'_0 \arrow[r,"i'"] \& A'.
\end{tikzcd}}
Sous ces conditions, le couple de foncteurs
\diagram{\begin{tikzcd}[ampersand replacement=\&]
 A_0 \arrow[r,shift left=2pt,"f_!^0"] \& A'_0 \arrow[l,shift left=2pt,"f_0^*"]
\end{tikzcd}}
est un couple de foncteurs adjoints. (NB quand je parle de couples de foncteurs adjoints, celui écrit au dessus de l'autre est l'adjoint à gauche.)

\uline{Dém} Soient
\[
 \deleted{\unsure{\xi}}\unsure{\eta|y}_0\in\deleted{\unsure{\mathrm{Ob}}}\mathrm{Ob}\,A_0,\quad \deleted{\unsure{\eta}}\xi'_0\in\mathrm{Ob}\,A'_0,
\]
il faut définir une bijection bifonctorielle
\[
 \mathrm{Hom}_{A'_0}(f_!^0(\xi'_0),\eta_0)\overset{?}{\simeq}
 \mathrm{Hom}_{A^{\unsure{\deleted{*}|'}}_0}(\xi'_0,f_0^*(\eta_0)).
\]
Pour cela, on va écrire $\unsure{\xi_0|\xi'_0}$ sous la forme
\[
 \xi'_0=\unsure{\rho|p}'(\xi')\qquad \xi'\in\mathrm{Ob}\,A',
\]
il suffira d'établir la bijection bifonctorielle en remplaçant $\xi'_0$ par $\rho'(\xi')$, i.e.
\[
 \mathrm{Hom}_{A^{\deleted{\unsure{\prime}}}_0}(f_!^0(\rho'(\xi')),\eta_0)\simeq
 \mathrm{Hom}_{A'_0}(\rho'(\xi'),f_0^{\unsure{d|*}}(\eta_0)).
\]
% archive sheet number: 203
